\part{Foundational Dynamics}
\section{Scope, System Boundary, and Non-Inference Note}

\noindent\textbf{What BaseCore does.}\enspace This document proves the mathematical base results that later layers may use: projection existence and uniqueness, strict contraction, unique fixed points, compact attractor families, parameterwise non-collapse under explicit anti-constant assumptions, inverse-limit identity under non-finite-determination assumptions, conditional rigidity, conditional non-simulability, and operational-criterion grammar.

\noindent\textbf{What BaseCore is not.}\enspace BaseCore is not the whole QICN corpus. It is not a bridge program, not runtime validation, not a phenomenality theory, not a human-equivalence claim, and not an external-adjudication layer.

\noindent\textbf{What should not be inferred.}\enspace Nothing in BaseCore alone licenses claims about consciousness, subjectivity, phenomenality, moral status, human-machine equivalence, bridge admissibility, or deployment readiness. Those are downstream burdens, not consequences of the theorems below.

\noindent\textbf{Source-of-truth rule.}\enspace The LaTeX tree is the source of truth. Compiled PDFs, release bundles, and runtime artifacts are derived surfaces only.

\section{Introduction}

\subsection{Purpose}

BaseCore is the canonical base layer of the Rigid Identity Framework. Its role is narrower and stronger than a broad architectural survey: it owns only the theorem-level mathematics that downstream papers may depend on without making BaseCore depend on them.

\subsection{Methodology}

We use standard tools from:
\begin{itemize}[leftmargin=*]
\item Hilbert-space geometry and metric projections;
\item contraction mappings and fixed-point theory;
\item compactness and continuity arguments;
\item projective systems and inverse limits;
\item assumption-led claim-boundary discipline.
\end{itemize}

\subsection{Scope discipline}

This section is pure mathematics. The symbols ``identity,'' ``regime,'' and ``criterion'' are structural labels only. Their ordinary-language interpretations are not part of any proof unless explicitly formalized as assumptions.

\section{Minimal Base Hypotheses}

\begin{hypothesis}[H1: Projection Structure]\label{hyp:H1}
Let $\Hilb$ be a real Hilbert space with inner product $\inner{\cdot}{\cdot}$ and norm $\norm{\cdot}$. Let $\I \subset \Hilb$ be non-empty, closed, and convex.
\end{hypothesis}

\begin{hypothesis}[H2: Contraction Operator]\label{hyp:H2}
Let $K:\Hilb\to\Hilb$ be a bounded linear operator with operator norm $\norm{K}<1$.
\end{hypothesis}

\begin{hypothesis}[H3: Completeness]\label{hyp:H3}
The ambient space $\Hilb$ is complete.
\end{hypothesis}

\begin{hypothesis}[H4: Parameter Compactness]\label{hyp:H4}
Let $\U$ be a compact metric space and let $\Gamma:\U\to\Hilb$ be uniformly continuous.
\end{hypothesis}

\begin{proposition}[Base sufficiency of H1--H4]\label{prop:minimal}
Hypotheses H1--H4 are sufficient for:
\begin{enumerate}[leftmargin=*]
\item existence and uniqueness of the metric projection onto $\I$;
\item strict contraction of the transition operators $T_u$;
\item existence and uniqueness of the fixed point $f_u^\ast$ for each $u\in\U$;
\item compactness of the attractor family $\A^\ast=\{f_u^\ast:u\in\U\}$.
\end{enumerate}
They do not, by themselves, imply parameterwise non-collapse.
\end{proposition}

\begin{proof}
The four items are exactly the content of Theorems~\ref{thm:projection}, \ref{thm:contraction}, \ref{thm:fixedpoint}, and \ref{thm:compactness}. Parameterwise non-collapse requires an additional anti-constant hypothesis introduced in Section~\ref{sec:parameterwise-noncollapse}.
\end{proof}

\section{The Projection Operator}

\begin{definition}[Metric Projection]\label{def:projection}
For $\Psi\in\Hilb$, define the metric projection onto $\I$ by
\[
\aleph(\Psi):=\arg\min_{\Phi\in\I}\norm{\Psi-\Phi}.
\]
\end{definition}

\begin{theorem}[Existence and Uniqueness of Projection]\label{thm:projection}
Under Hypothesis~\ref{hyp:H1}, the projection $\aleph(\Psi)$ exists and is unique for every $\Psi\in\Hilb$.
\end{theorem}

\begin{proof}
Existence follows because a minimizing sequence in the closed convex set $\I$ is Cauchy by the parallelogram law, hence converges in the complete space $\Hilb$ to a point of $\I$ that attains the infimum. Uniqueness follows from strict convexity of the norm induced by the Hilbert inner product: if two minimizers existed, their midpoint would produce the same minimum and the parallelogram identity would force the two minimizers to coincide.
\end{proof}

\begin{lemma}[Non-Expansiveness of the Projection]\label{lem:nonexp}
The map $\aleph:\Hilb\to\I$ is non-expansive:
\[
\norm{\aleph(\Psi_1)-\aleph(\Psi_2)}\le \norm{\Psi_1-\Psi_2}
\qquad
\forall \Psi_1,\Psi_2\in\Hilb.
\]
\end{lemma}

\begin{proof}
This is the standard Hilbert-space projection inequality. Let $\Phi_i=\aleph(\Psi_i)$. The projection characterization
\[
\inner{\Psi_i-\Phi_i}{\Phi-\Phi_i}\le 0
\qquad
\forall \Phi\in\I
\]
applied with $\Phi=\Phi_2$ and $\Phi=\Phi_1$ yields
\[
\norm{\Phi_1-\Phi_2}^2
\le
\inner{\Psi_1-\Psi_2}{\Phi_1-\Phi_2}
\le
\norm{\Psi_1-\Psi_2}\norm{\Phi_1-\Phi_2}.
\]
If $\Phi_1\neq \Phi_2$, divide by $\norm{\Phi_1-\Phi_2}$.
\end{proof}

\section{The Transition Operator}

\begin{definition}[Transition Operator]\label{def:transition}
For each $u\in\U$, define
\[
T_u(\Psi):=\aleph(K\Psi+\Gamma(u)).
\]
\end{definition}

\begin{theorem}[Contractivity]\label{thm:contraction}
Under Hypotheses~\ref{hyp:H1} and \ref{hyp:H2}, each $T_u$ is a strict contraction with Lipschitz constant at most $\norm{K}$:
\[
\norm{T_u(\Psi_1)-T_u(\Psi_2)}\le \norm{K}\,\norm{\Psi_1-\Psi_2}.
\]
\end{theorem}

\begin{proof}
By Lemma~\ref{lem:nonexp},
\begin{align*}
\norm{T_u(\Psi_1)-T_u(\Psi_2)}
&=
\norm{\aleph(K\Psi_1+\Gamma(u))-\aleph(K\Psi_2+\Gamma(u))} \\
&\le \norm{K(\Psi_1-\Psi_2)}
\le \norm{K}\,\norm{\Psi_1-\Psi_2}.
\end{align*}
Since $\norm{K}<1$, strict contraction follows.
\end{proof}

\section{Existence and Compactness of Attractors}

\begin{theorem}[Unique Fixed Point]\label{thm:fixedpoint}
Under Hypotheses~\ref{hyp:H1}--\ref{hyp:H3}, for each $u\in\U$ there exists a unique $f_u^\ast\in\Hilb$ such that
\[
T_u(f_u^\ast)=f_u^\ast.
\]
Moreover, for every $\Psi_0\in\Hilb$,
\[
\lim_{n\to\infty}T_u^n(\Psi_0)=f_u^\ast.
\]
\end{theorem}

\begin{proof}
By Theorem~\ref{thm:contraction}, each $T_u$ is a strict contraction on the complete metric space $(\Hilb,\norm{\cdot})$. The Banach fixed-point theorem yields existence, uniqueness, and convergence of iterates.
\end{proof}

\begin{definition}[Attractor Family]\label{def:attractor}
Define the attractor family
\[
\A^\ast:=\{f_u^\ast:u\in\U\}.
\]
\end{definition}

\begin{theorem}[Compactness of the Attractor Family]\label{thm:compactness}
Under Hypotheses~\ref{hyp:H1}--\ref{hyp:H4}, the set $\A^\ast$ is compact in $\Hilb$.
\end{theorem}

\begin{proof}
Let $F:\U\to\Hilb$ be given by $F(u)=f_u^\ast$. The fixed points depend continuously on $u$ because the contraction constant is uniform and $\Gamma$ is uniformly continuous on the compact parameter space. Since $\U$ is compact and $\A^\ast=F(\U)$, the image is compact.
\end{proof}

\section{Parameterwise Non-Collapsibility}\label{sec:parameterwise-noncollapse}

\begin{definition}[Constant Subspace]
Whenever the ambient realization distinguishes constant elements, let $\mathcal{N}\subset\Hilb$ denote the closed subspace of constants.
\end{definition}

\begin{definition}[Parameterwise Non-Collapsible Family]
A family $\{T_u\}_{u\in\U}$ is parameterwise non-collapsible if
\[
f_u^\ast\notin\mathcal{N}
\qquad
\forall u\in\U.
\]
\end{definition}

\begin{hypothesis}[H5: No Constant Fixed Points]\label{hyp:H5}
For every parameter $u\in\U$ and every constant element $c\in\mathcal{N}$,
\[
T_u(c)\neq c.
\]
\end{hypothesis}

\begin{theorem}[Structural Non-Collapsibility]\label{thm:noncollapse}
Under Hypotheses~\ref{hyp:H1}--\ref{hyp:H5}, the family is parameterwise non-collapsible:
\[
f_u^\ast\notin\mathcal{N}
\qquad
\forall u\in\U.
\]
\end{theorem}

\begin{proof}
Fix $u\in\U$. Suppose by contradiction that $f_u^\ast\in\mathcal{N}$. Write $f_u^\ast=c$ with $c\in\mathcal{N}$. Since $f_u^\ast$ is a fixed point,
\[
T_u(c)=T_u(f_u^\ast)=f_u^\ast=c,
\]
contradicting Hypothesis~\ref{hyp:H5}. Therefore $f_u^\ast\notin\mathcal{N}$.
\end{proof}

\begin{remark}[What H5 does not say]
Hypothesis~\ref{hyp:H5} is parameterwise. It does not merely say that \emph{some} parameter breaks constant collapse. It excludes constant fixed points for \emph{every} parameter. This is exactly the quantifier strength needed for Theorem~\ref{thm:noncollapse}.
\end{remark}

\section{Summary of Base Results}

\begin{enumerate}[leftmargin=*]
\item H1--H4 give a complete theorem-tight base for projection existence, strict contraction, unique fixed points, and compact attractor families.
\item H5 is a separate anti-constant assumption block used only for parameterwise non-collapse.
\item No empirical, phenomenological, biological, or runtime claim is inferred at this level.
\item Later sections may add conditional identity, rigidity, and criterion theorems, but only by introducing their own explicit assumption blocks.
\end{enumerate}
